% !TeX root = ../papers/nc-p3.tex
\section{反対階数15の二つの四重被覆とその分岐曲線}
\label{sec:rank-fifteen-covers}
\begin{theorem}\label{thm:rank-fifteen-covers}
正の固定曲線の通常倍率平面上で
\[
 k=\frac{(\lambda+1)(\lambda z+1)}{4\lambda z},\qquad
 m=-\frac{(\lambda-1)(\lambda z-1)}{4\lambda z},\qquad
 P(n)=n^4+(1-zk^2)n^2+zm^2
\]
と置く。反対階数が低下する全窓は、斉次座標 $[a:b:e:f]$ で次の二表示に尽きる。
各表示で $P(n)=0$ とし、
\[
 (a,b,e,f)=(1+n,k+m/n,1-n,k-m/n)
\]
を取り、第二表示ではこれを $(zb,a,zf,e)$ に移す。
二表示はそれぞれ有限平坦な四重被覆を与え、幾何学的一般モノドロミーは位数8の二面体群である。
二位置移動は各被覆で $n\mapsto-n$ である。
\begin{enumerate}
\item $zk^2=1$ では二表示が一致し、四点全てが前定理の正則な階数14窓である。
\item $\delta=(1-zk^2)^2-4zm^2=0$ では合計四つの相異なる階数15窓があり、各根の重複度は2である。
\item その他の通常倍率では合計八つの相異なる階数15窓がある。
\end{enumerate}
個数は固定した座標での幾何パラメータ数である。
\end{theorem}
\begin{proof}
$C=\lambda^2z+1$ と書く。前の交換行列による $(e,f)=u(N,D)$ において
$af+be=C u(a^2-zb^2)$、$ae+zbf=2\lambda z u(a^2-zb^2)$ である。
補完の二つの二次式の残差はともに
$G=\lambda(z+1)u(a^2-zb^2)-ab-ef$ の倍数で、分子の倍率はそれぞれ $D,N$、
共通分母は $\lambda(z+1)(a^2-zb^2)$ となる。
$i$ 番目の二次式の選択が対根に対応することは $zf_1 f+e_1 e=0$ から従う。
従って許容座標では補完と交換の二条件が次の二次式に一致する。
\[
 2\lambda z(af+be)=C(ae+zbf),\qquad
 2\lambda z(ab+ef)=\lambda(z+1)(ae+zbf).
\]
逆にも第一式から $[e:f]=[N:D]$ が従い、第二式で $G=0$ となる。$C=0$ の場合も
交換行列は可逆であり、この議論は $C$ を分母に使わない。

最初の階数低下条件 $a^2+e^2=2zbf$ では $2(ae+zbf)=(a+e)^2$ である。
$a+e=0$ なら二式の和差から $(a-e)(b-f)=0$ となり、一方の行列が零になるか
$a^2=zb^2$ となるので除かれる。$a+e=1$ と規格化すると $b+f=k$、
$(a-e)(b-f)=m$ であり、$n=a-e$ として残条件が $P(n)=0$ になる。
もう一つの階数低下条件は上記の座標交換で移る。この交換は二つの補完式も保つ。
得られた全根で $n\ne0$ であり、$a^2-zb^2,e^2-zf^2,a^2-zf^2,e^2-zb^2$ は非零である。
\path{computations/verify_rank_fifteen_covers.py} は分母を払ったこれらの式と $P$ の
終結式が全て通常倍率の単元であることを両表示で検算する。このため座標端でも許容根を失わない。
第一表示で $a+e=1$ とした他方の階数低下条件の左辺 $z(b^2+f^2)-2ae$ は $zk^2-1$ である。
根数は $n^2$ の二次式の判別式と $zm^2\ne0$ から従う。

最後に $L=\mathbb C(z,\lambda)$ 上、$zm^2$ は平方でない。
$\delta$ も $z\delta$ も平方でない。実際 $\lambda=2$ で $\delta$ の既約分数の
分子は正次数の平方自由式で $z=0$ に根を持たず、その分母は平方である。
$n^2$ の二根を $v_1,v_2$ とすると、その二次分解体 $L(\sqrt\delta)$ においても
$v_1v_2=zm^2$ は非平方である。共役な $v_1,v_2$ は従って両方非平方で、
平方類も独立である。分解体の次数は8であり、二組の符号と二根の交換から二面体群を得る。
\end{proof}

\begin{proposition}\label{prop:rank-fifteen-ramification-curve}
各表示の分岐上の被約な窓の曲線は幾何学的に既約で、その滑らかな完備化は
$y^2=-t^4+14t^2-1$、$j=35152/9$ の種数1曲線である。
両曲線上で一般窓は標準Koszul左右AS正則である。
\end{proposition}
\begin{proof}
分岐上で $1-zk^2=2rm$ により $r^2=z$ の符号を一意に取り、
$w=\lambda r$、$K=rk$ と置く。$rk,rm$ は
$r+1/r$ と $w+1/w$ の和・差の四分の一である。従って
\[
 r+1/r=2-(K-1)^2,\qquad
 w+1/w=(K+1)^2-2,\qquad n^2=(K^2-1)/2.
\]
$K=\pm1$ は $m=0$ を与えて除かれる。
$A=(2r-2+(K-1)^2)/(K-1)$、$B=(2w+2-(K+1)^2)/(K+1)$ と置くと
$A^2=(K-3)(K+1)$、$B^2=(K-1)(K+3)$ となる。
$U=A/(K+1),V=n/(K+1)$ は $4V^2-U^2=1$ を満たす。
$t=2V+U$ を用いて
\[
 U=(t-t^{-1})/2,\quad V=(t+t^{-1})/4,\quad
 K=4/(1-U^2)-1,\quad y=2tB/n
\]
とすれば所定の四次式が従う。これらの式は両方向の有理逆を与える。
四根は相異なり、二元四次式の不変量から記載した $j$ を得る。

$K=2$ の八つの $(r,w,n)$ は139で全て単純分解する。二表示を合わせた16窓について、
上記コードはテンソルの隣接階数6、反対階数15、経路次元35,56、中心の文字次元
$(3,2,0,0)$ を調べ、全四始点の四中心商で長さ13まで標準次元を得る。
正の固定族の一般点からの経路・中心次元の下界と命題\ref{prop:fixed-residual-regularity}
の持上げにより、この16幾何点の全次数の正則性が従う。
定階数の中心から四式を近傍に延ばせるため、各既約分岐曲線でも非空開集合が正則となる。
\end{proof}
